\documentclass{article}


\usepackage[style=authortitle]{biblatex} % bibliografhy
\usepackage{graphicx}

% added tikz, and a mysterius package... just to try out.
\usepackage{tikz}
\usepackage{tikz,fullpage}
\usetikzlibrary{arrows,%
                petri,%
                topaths}%
\usepackage{tkz-berge}

\addbibresource{./tex/INF221.bib}

\begin{document}

\begin{titlepage}
	\begin{center}
		\vspace*{1cm}
		\Huge
		\textbf{Analyses of various algorithms}\\

		\vspace{0.5cm}
		\LARGE

		Term paper for INF221 2021\\

		\vspace{1.5cm}

		\textit{Mats Hoem Olsen, Roy Erling Granheim, Sindre Stokke}\\

		\vfill

		\includegraphics[width=0.4\textwidth]{./tex/pic/logo.eps}\\
		
		\LARGE
		DataSci\\
		Norwegian university of life science\\
		2021 December 3
	\end{center}		
\end{titlepage}

\tableofcontents
\begin{abstract}
algorithm goes brrrr.
\end{abstract}
\newpage

\section{Introduction}

Hi, we are students. How are you?\parencite[see][]{cormen_introduction_2009}

\section{Theory}


Provide a brief description of the algorithms you will be
investigating, including pseudocode for the algorithms. Describe in
particular the expected runtime of algorithms in terms of problem
size.  Use a separate subsection for each algorithm.

\subsection{Algo 1: Bubble sort}\label{sec:algo1}

\begin{algorithm}[H]
  \caption{Insertion sort algorithm from \cite[Ch.~2.1]{cormen_introduction_2009}.}
  \label{lst:bubble_sort}
  \begin{codebox}
    \Procname{$\proc{Bubblesort}(A)$}
    \li \For $i \gets 1$ \To $\attrib{A}{length} - 1$
    \li \Do \For $j \gets \attrib{A}{length}$ \Downto $i + 1$
    \li 	\Do \If $A[j] < A[j - 1]$
    \li         \Then exchange $A[j]$ with $A[j - 1]$
    \End    
    \End
    \End
  \end{codebox}
\end{algorithm}

Pseudocode for the first algorithm is shown in
Listing~\ref{lst:bubble_sort}. Best case runtime for this algorithm
is

\begin{equation}
  T(n) = \Theta(n) \;.  \label{eq:ins_sort_best}
\end{equation}

It is achieved for correctly sorted input data.

\subsection{Algo 2: Insertion sort}\label{sec:algo2}

\begin{algorithm}[H]
  \caption{Insertion sort algorithm from \cite[Ch.~2.1]{cormen_introduction_2009}.}
  \label{lst:insertion_algo}
  \begin{codebox}
    \Procname{$\proc{Insertion-Sort}(A)$}
    \li \For $j \gets 2$ \To $\attrib{A}{length}$
    \li \Do
    $\id{key} \gets A[j]$
    \li     $i \gets j-1$
    \li      \While $i>0$ and $A[i] > \id{key}$
    \li      \Do
    $A[i+1] \gets A[i]$
    \li         $i \gets i-1$
    \End    
    \li       $A[i+1]\gets \id{key}$
    \End
  \end{codebox}
\end{algorithm}

Pseudocode for the second algorithm is shown in
Listing~\ref{lst:insertion_algo}. Best case runtime for this algorithm
is

\begin{equation}
  T(n) = \Theta(n) \;.  \label{eq:ins_sort_best}
\end{equation}

It is achieved for correctly sorted input data.

\subsection{Algo 3: Quicksort}\label{sec:algo3}

\begin{algorithm}[H]
  \caption{Quicksort algorithm from \cite[Ch.~2.1]{cormen_introduction_2009}.}
  \label{lst:quicksort_algo}
  \begin{codebox}
    \Procname{$\proc{Quicksort}(A, p, r)$}
    \li \If $p < r$
    \li \Then $q \gets \proc{Partition}(A, p, r)$
    \li     $\proc{Quicksort}(A, p, q - 1)$
    \li     $\proc{Quicksort}(A, q + 1, r)$
    \End
  \end{codebox}
\end{algorithm}

\begin{algorithm}[H]
  \caption{Partition from \cite[Ch.~2.1]{cormen_introduction_2009}.}
  \label{lst:partition}
  \begin{codebox}
    \Procname{$\proc{Partition}(A, p, r)$}
    \li $x \gets A[r]$
    \li $i = p - 1$
    \li \For $j \gets p$ \To $r - 1$
    \li \Do \If $A[j] \le x$
    \li  	\Then $i \gets i + 1$
    \li     exchange $A[i]$ with $A[j]$
    \End
    \End
    \li exchange $A[i + 1]$ with $A[r]$
    \li \Return $i + 1$
  \end{codebox}
\end{algorithm}

Pseudocode for the second algorithm is shown in
Listing~\ref{lst:partition}. Best case runtime for this algorithm
is

\begin{equation}
  T(n) = \Theta(n) \;.  \label{eq:ins_sort_best}
\end{equation}

It is achieved for correctly sorted input data.

\subsection{Algo 4: Quicksort Median of Three}\label{sec:algo4}

\begin{algorithm}[H]
  \caption{Quicksort median of three algorithm from \cite[Ch.~2.1]{cormen_introduction_2009}.}
  \label{lst:quickmed3_algo}
  \begin{codebox}
    \Procname{$\proc{Insertion-Sort}(A)$}
    \li \For $j \gets 2$ \To $\attrib{A}{length}$
    \li \Do
    $\id{key} \gets A[j]$
    \li     $i \gets j-1$
    \li      \While $i>0$ and $A[i] > \id{key}$
    \li      \Do
    $A[i+1] \gets A[i]$
    \li         $i \gets i-1$
    \End    
    \li       $A[i+1]\gets \id{key}$
    \End
  \end{codebox}
\end{algorithm}

Pseudocode for the second algorithm is shown in
Listing~\ref{lst:quickmed3_algo}. Best case runtime for this algorithm
is

\begin{equation}
  T(n) = \Theta(n) \;.  \label{eq:ins_sort_best}
\end{equation}

It is achieved for correctly sorted input data.

\subsection{Algo 5: Quicksort Insertion sort hybrid}\label{sec:algo5}

\begin{algorithm}[H]
  \caption{Insertion sort algorithm from \cite[Ch.~2.1]{cormen_introduction_2009}.}
  \label{lst:quickinsert_algo}
  \begin{codebox}
    \Procname{$\proc{Insertion-Sort}(A)$}
    \li \For $j \gets 2$ \To $\attrib{A}{length}$
    \li \Do
    $\id{key} \gets A[j]$
    \li     $i \gets j-1$
    \li      \While $i>0$ and $A[i] > \id{key}$
    \li      \Do
    $A[i+1] \gets A[i]$
    \li         $i \gets i-1$
    \End    
    \li       $A[i+1]\gets \id{key}$
    \End
  \end{codebox}
\end{algorithm}

Pseudocode for the second algorithm is shown in
Listing~\ref{lst:quickinsert_algo}. Best case runtime for this algorithm
is

\begin{equation}
  T(n) = \Theta(n) \;.  \label{eq:ins_sort_best}
\end{equation}

It is achieved for correctly sorted input data.

\subsection{Algo 6: Mergesort}\label{sec:algo6}
why are you like this

\begin{algorithm}[H]
  \caption{Mergesort algorithm from \cite[Ch.~2.1]{cormen_introduction_2009}.}
  \label{lst:mergesort_algo}
  \begin{codebox}
    \Procname{$\proc{Merge-Sort}(A, p, r)$}
	\li \If $p < r$
	\li \Then $q \gets \left\lfloor(p + r) / 2\right\rfloor$
	\li 	$\proc{Merge-Sort}(A, p, q)$
	\li 	$\proc{Merge-Sort}(A, q+1, r)$
	\li 	$\proc{Merge}(A, p, q, r)$
	\End
  \end{codebox}
\end{algorithm}

\begin{algorithm}[H]
  \caption{Merge from \cite[Ch.~2.1]{cormen_introduction_2009}.}
  \label{lst:merge}
  \begin{codebox}
    \Procname{$\proc{Merge}(A, p, q, r)$}
	\li $n_1 \gets q-p+1$
	\li $n_2 \gets r-q$
	\li let $L[1..n_1+1]$ and $R[1..n_2+1]$ be new arrays
	\li \For $i \gets 1$ \To $n_1$
	\li \Do $L[i] \gets A[p+i-1]$
	\End
	\li \For $j \gets 1$ \To $n_2$
	\li \Do $R[j] \gets A[q+j]$
	\End
	\li $L[n_1+1] \gets \infty$
	\li $R[n_2+1] \gets \infty$
	\li $i \gets 1$
	\li $j \gets 1$
	\li \For $k \gets p$ \To $r$
	\li \Do \If $L[i] \le R[j]$
	\li 	\Then $A[k] \gets L[i]$
	\li 		$i \gets i+1$
	\li 	\Else $A[k] \gets R[j]$
	\li 	$j \gets j+1$
	\End
	\End
  \end{codebox}
\end{algorithm}

Pseudocode for the second algorithm is shown in
Listing~\ref{lst:merge}. Best case runtime for this algorithm
is

\begin{equation}
  T(n) = \Theta(n) \;.  \label{eq:ins_sort_best}
\end{equation}

It is achieved for correctly sorted input data.

\subsection{Algo 7: Mergesort Insertion sort hybrid}\label{sec:algo7}

\begin{algorithm}[H]
  \caption{Insertion sort algorithm from \cite[Ch.~2.1]{cormen_introduction_2009}.}
  \label{lst:mergeinsert_algo}
  \begin{codebox}
    \Procname{$\proc{Insertion-Sort}(A)$}
    \li \For $j \gets 2$ \To $\attrib{A}{length}$
    \li \Do
    $\id{key} \gets A[j]$
    \li     $i \gets j-1$
    \li      \While $i>0$ and $A[i] > \id{key}$
    \li      \Do
    $A[i+1] \gets A[i]$
    \li         $i \gets i-1$
    \End    
    \li       $A[i+1]\gets \id{key}$
    \End
  \end{codebox}
\end{algorithm}

Pseudocode for the second algorithm is shown in
Listing~\ref{lst:mergeinsert_algo}. Best case runtime for this algorithm
is

\begin{equation}
  T(n) = \Theta(n) \;.  \label{eq:ins_sort_best}
\end{equation}

It is achieved for correctly sorted input data.

\subsection{Algo 8: Python sort}\label{sec:algo8}

Best case runtime for this algorithm is

\begin{equation}
  T(n) = \Theta(n) \;.  \label{eq:ins_sort_best}
\end{equation}

It is achieved for correctly sorted input data.

\subsection{Algo 9: Numpy sort}\label{sec:algo9}

Best case runtime for this algorithm is

\begin{equation}
  T(n) = \Theta(n) \;.  \label{eq:ins_sort_best}
\end{equation}

It is achieved for correctly sorted input data.


\section{Method used}

In this section we will show our workflow and automation of tests.\parencite[look at ][time: 1:23]{brady_haran_death_2020}
\begin{figure}
	\centering
	\includegraphics[]{./pics/flower.pdf}
\end{figure}

\section{Benchmarks}

In this section we will analyse the performance of the different algorithms on exponentially growing test data.

\section{Discussion}

We will in this section discuss our findings, and thoughts\footnote{Although these thoughts are subjective, we will present them as facts. For our thoughts feel real to us.}

\section{Acknowledgements}

We acknowledge that we may, or may not be made out of bread.

\printbibliography
\end{document}
